\documentclass[french,12pt]{report}
%Some packages I commonly use.
%\usepackage[english]{babel}
\usepackage[french]{babel}
\usepackage{graphicx}
\usepackage{framed}
\usepackage[normalem]{ulem}
\usepackage{amsmath}
\usepackage{amsthm}
\usepackage{amssymb}
\usepackage{amsfonts}
\usepackage{enumerate}
\usepackage{xcolor}
\usepackage[utf8]{inputenc}
\usepackage[top=1 in,bottom=1in, left=1 in, right=1 in]{geometry}
\usepackage{systeme}
\usepackage{hyperref} 
\usepackage{dsfont}
\usepackage[
backend=biber,
style=alphabetic,
sorting=ynt
]{biblatex}
 
%A bunch of definitions that make my life easier
\newcommand{\matlab}{{\sc Matlab} }
\newcommand{\cvec}[1]{{\mathbf #1}}
\newcommand{\rvec}[1]{\vec{\mathbf #1}}
\newcommand{\ihat}{\hat{\textbf{\i}}}
\newcommand{\jhat}{\hat{\textbf{\j}}}
\newcommand{\khat}{\hat{\textbf{k}}}
\newcommand{\minor}{{\rm minor}}
\newcommand{\trace}{{\rm trace}}
\newcommand{\spn}{{\rm Span}}
\newcommand{\rem}{{\rm rem}}
\newcommand{\ran}{{\rm range}}
\newcommand{\range}{{\rm range}}
\newcommand{\mdiv}{{\rm div}}
\newcommand{\proj}{{\rm proj}}
\newcommand{\R}{\mathbb{R}}
\newcommand{\N}{\mathbb{N}}
\newcommand{\Q}{\mathbb{Q}}
\newcommand{\Z}{\mathbb{Z}}
\newcommand{\<}{\langle}
\renewcommand{\>}{\rangle}
\renewcommand{\emptyset}{\varnothing}
\newcommand{\attn}[1]{\textbf{#1}}
\theoremstyle{definition}
\newtheorem*{remark}{Remarque}
\newtheorem{theorem}{Theorem}
\newtheorem{corollary}{Corollary}
\newtheorem*{definition}{Definition}
\newtheorem*{example}{Example}
\newtheorem*{note}{Note}
\newtheorem{exercise}{Exercise}
\newcommand{\bproof}{\bigskip {\bf Proof. }}
\newcommand{\eproof}{\hfill\qedsymbol}
\newcommand{\Disp}{\displaystyle}

\setlength{\columnseprule}{1 pt}
\frenchsetup{StandardLists=true}

\addbibresource{biblio.bib}



\begin{document}

\begin{titlepage}

\begin{figure}
     \begin{minipage}{0.4\textwidth}
         \begin{flushleft}
             \includegraphics[scale=0.7]{univ.png}
             
         \end{flushleft}
     \end{minipage}
     \hfill
     \begin{minipage}{0.4\textwidth}
         \begin{flushright}
             \includegraphics[scale=0.3]{IAE.png}
          
         \end{flushright}
     \end{minipage}
\end{figure}


\vspace{4cm}

\center
\textsc{\LARGE Master 2 MIND/SIAD}\\
\vspace{1cm}
\textrm{\textbf{Statistique Computationnelle}}

\rule{\linewidth}{1pt}\\[0.4cm]

\huge {Projet 4 : Inférence bayésienne pour un modèle log-linéaire }
\rule{\linewidth}{1pt}\\[0.5cm]

 
 




\vspace{1cm}

\Large Proposé par : \\[0.5cm]

\LARGE \textbf{Jean-Michel Marin}\\






     




 
\vspace{1cm}

\Large Fait par : \\[0.9cm]


\LARGE \textbf{Mégane} \textsc{ \textbf{Dieval}}\\  
\vspace{0.5cm}
\LARGE \textbf{Ibrahim} \textsc{ \textbf{Gaizi}}\\
\vspace{0.5cm}
\LARGE \textbf{Oumayma} \textsc{ \textbf{Khalifi}}\\     


  

\centering




     




\end{titlepage}






\textbf{\textit{Question 1 : }} \\ \\

On a : $$ log(\mu_i)=x_i^T\beta$$

Or nous avons $$l(\mu/y)=\prod_{i=1}^n \frac{1}{y_i}(\mu_i^{y_i}exp(-\mu_i)) $$ 
et on a $$ log(\mu_i)=x_i^T\beta$$
On applique maintenant la formule de Bayes: 
$$\pi(\beta,X) = \pi(\beta/X)\pi(X)$$
$$\pi(\beta/y) \propto l(\beta/y)\pi(\beta) $$
$$\propto \prod_{i=1}^n \frac{1}{y_i} \exp(y_i X_i^T\beta)\exp(-\exp(X_i^T\beta))\frac{\sqrt{n}}{\sqrt{2\pi}} (X¨TX)^{\frac{1}{2}} \exp(-\frac{1}{2}(X_i^TX^TXX_i) $$
$$\propto (X^TX)^{\frac{1}{2}}\exp(\sum_{i=1}^n y_i X_i^T\beta-\exp(\x_i^T\beta-\frac{1}{2}X_i^T(X^TX)X_i)$$



\end{document}




